\mainchapter{Marco Teórico}{\topbanner}\label{cap:marco_teorico}
\vspace*{0.7cm}

En este capítulo se profundiza en los fundamentos teóricos que soportan el desarrollo del proyecto, complementando los conceptos presentados en el capítulo anterior. Se describen los modelos y estándares asociados a la gestión de identidades y al control de acceso, los mecanismos de autenticación definidos por el protocolo \acs{LDAP}, así como las normas de calidad y seguridad que orientan las decisiones adoptadas.
\section{Modelos de gestión de identidades}

La gestión de identidades se sustenta en marcos conceptuales y normativos que definen los elementos involucrados en la representación digital de una entidad. La norma ISO/IEC 24760-1 define una identidad como la representación de una entidad dentro de un contexto específico, compuesta por un conjunto de atributos que la caracterizan y un identificador que permite distinguirla de otras, e identifica tres fases fundamentales en su administración: el establecimiento de la identidad, su mantenimiento durante el ciclo de vida y su finalización \citep{ISO24760}. Este ciclo de vida comprende actividades como el registro inicial de la identidad, la actualización de sus atributos, la suspensión temporal y el retiro definitivo de las cuentas cuando dejan de ser necesarias.

Por otra parte, \citet{Pohn2022} proponen un modelo de referencia para la gestión de identidades que organiza sus funcionalidades en torno a procesos como la provisión de cuentas, la autenticación, la autorización, la federación y la auditoría. De acuerdo con este modelo, una arquitectura de gestión de identidades debe garantizar la consistencia de los datos de identidad entre los distintos sistemas que los consumen, lo cual constituye uno de los principales argumentos a favor de la centralización mediante servicios de directorio. Esta visión coincide con las limitaciones identificadas en los entornos donde cada aplicación administra sus propios repositorios de usuarios, situación que incrementa la carga operativa y favorece la aparición de inconsistencias \citep{Pohn2020}.

\section{Modelos de control de acceso}

Los modelos de control de acceso establecen las reglas mediante las cuales se determinan los privilegios de las entidades sobre los recursos de un sistema. Dentro de estos modelos, el control de acceso discrecional delega la decisión sobre los permisos al propietario de cada recurso, mientras que el control de acceso obligatorio asigna dichos permisos mediante políticas centrales administradas por el sistema \citep{RFC4949}.

Un tercer modelo ampliamente adoptado es el control de acceso basado en roles (\acs{RBAC}), formalizado por \citet{Sandhu1996}. En este modelo, los permisos no se asignan directamente a los usuarios, sino a roles que agrupan las funciones desempeñadas dentro de la organización; de esta manera, un usuario adquiere los privilegios correspondientes al pertenecer a un rol determinado. Este enfoque simplifica la administración en entornos con gran cantidad de usuarios, facilita la aplicación del principio de mínimo privilegio y refleja de manera natural la estructura organizacional. Los grupos y roles disponibles en los servicios de directorio constituyen la materialización práctica de este modelo, permitiendo gestionar los accesos de forma agregada y coherente con la estructura institucional.

\section{Autenticación en \texorpdfstring{\acs{LDAP}}{LDAP}}

El protocolo \acs{LDAP} define en su especificación los mecanismos mediante los cuales una entidad establece una sesión autenticada con el servicio de directorio. La operación de vinculación admite dos grandes categorías de autenticación: la autenticación simple, en la cual se presenta un nombre distintivo y una contraseña, y la autenticación mediante el nivel de seguridad simple y capa de protección (\acs{SASL}), que permite incorporar mecanismos externos de autenticación como Kerberos \citep{RFC4513}.

Dado que la autenticación simple transmite las credenciales a través de la red, es indispensable proteger el canal de comunicación. Para ello, \acs{LDAP} contempla el uso de \acs{TLS}, ya sea mediante conexiones dedicadas cifradas (\acs{LDAPS}) o mediante el mecanismo STARTTLS, que eleva una conexión existente a una protegida mediante cifrado \citep{RFC4513}. La protección del canal garantiza la confidencialidad e integridad de las credenciales y de la información intercambiada entre los clientes y el servicio de directorio, constituyendo un requisito básico para el despliegue seguro del servicio.

\section{Normas de calidad y seguridad aplicables}

El diseño e implementación de un servicio de directorio se orienta según normas internacionales que establecen buenas prácticas de seguridad y criterios de calidad. La norma ISO/IEC 27000 proporciona el vocabulario y los conceptos generales de la seguridad de la información, definiendo la preservación de la confidencialidad, la integridad y la disponibilidad como objetivos fundamentales \citep{ISO27000}. Sobre esta base, la norma ISO/IEC 27001 establece requisitos para un sistema de gestión de seguridad de la información, incluyendo controles específicos relacionados con la gestión de accesos, la gestión de identidades, la autenticación y la distribución de privilegios \citep{ISO27001}, los cuales resultan directamente aplicables al servicio de directorio propuesto.

En cuanto a los mecanismos de autenticación, las directrices de identidad digital del \textit{National Institute of Standards and Technology} definen niveles de garantía para la identidad, la autenticación y la federación, recomendando la autenticación multifactor como medida para elevar la resistencia frente a ataques dirigidos al robo de credenciales \citep{NIST80063v4}. Por su parte, la norma ISO/IEC 25010 aporta un modelo de calidad para sistemas software, cuyas características de seguridad, fiabilidad, compatibilidad, mantenibilidad y adecuación funcional permiten evaluar de manera integral la solución implementada \citep{ISO25010}. En conjunto, estas normas proporcionan el referente normativo que guía las decisiones de diseño, implementación y verificación del presente proyecto.

\section{\texorpdfstring{\acs{PMBOK}}{PMBOK}}

La dirección de proyectos cuenta con marcos de referencia estandarizados que consolidan las buenas prácticas de la disciplina. Entre ellos, la guía \acs{PMBOK}, publicada por el \textit{Project Management Institute} (\acs{PMI}), es reconocida internacionalmente como un estándar para la dirección de proyectos, al reunir procesos, herramientas, técnicas y principios aplicables a proyectos de diversa naturaleza \citep{PMI2021}. El estándar asociado a esta guía estructura el ciclo de vida de un proyecto en cinco grupos de procesos: inicio, planificación, ejecución, seguimiento y control, y cierre, los cuales delimitan las actividades realizadas desde la autorización del proyecto hasta su finalización \citep{PMI2021}. En el presente proyecto, esta guía se adopta como referente para orientar la planificación, la organización de las actividades y el seguimiento de las fases desarrolladas.
